The basic data set underlying our empirical work is a balanced panel of 50 states over 26 quarters from 2019-Q1 to 2025-Q2. We examine three main groups of outcomes: GLP-1 prescription drug utilization, gross Medicaid spending on GLP-1 prescriptions, and net-of-rebate Medicaid spending on GLP-1 prescriptions. All three outcomes are computed separately for obesity-indicated and diabetes/cardiovascular-indicated GLP-1 drugs and are measured at the state-quarter level. We use State Drug Utilization Data (SDUD) to measure total prescriptions and total gross spending for each group of drugs by state and quarter. To estimate per-unit rebates for each GLP-1 NDC code, we use drug pricing data from CMS National Average Drug Acquisition Cost (NADAC) and VA Federal Supply Schedule (FSS) as proxies for the confidential elements of the MDRP rebate formula.  The per-unit rebates are combined with SDUD data on total Medicaid financed units of each NDCs to compute the basic and inflationary rebates, which are then aggregated to used to compute total net-of-rebate Medicaid spending for each group of drugs by state and quarter. We scale the prescription, spending, and net-spending totals by the Medicaid enrollee-months, which we compute using the monthly enrollment data in CMS Medicaid and CHIP Performance Indicators (CMSPI). These enrollee-month counts are also used to construct enrollment weights in our analysis.

\paragraph{State Drug Utilization Data.} The SDUD reports total prescriptions, units dispensed, and Medicaid reimbursements by state, plan type (MCO vs FFS), and quarter for every NDC code that had any Medicaid prescriptions during the quarter (2019-Q1 to 2025-Q2). We restricted the SDUD sample to NDC codes for GLP-1 medications, classifying each NDC as an obesity-indicated (Saxenda, Wegovy, and Zepbound) or diabetes/cardiovascular-indicated formulation (Victoza, Ozempic, Rybelsus, Mounjaro, exenatide, dulaglutide, and lixisenatide).\footnote{We identified NDC codes for GLP-1 medications using the National Library of Medicine (NLM)'s RxMix database. Our list includes all NDCs for liraglutide (Saxenda, Victoza), semaglutide (Wegovy, Ozempic, Rybelsus), tirzepatide (Zepbound, Mounjaro), and other GLP-1 receptor agonists (exenatide, dulaglutide, lixisenatide). The complete list of NDCs is available in the replication archive.} The SDUD suppresses information on the number of prescriptions for cells with fewer than 11 prescriptions. We take the midpoint and assume that each suppressed cell had 5 prescriptions. We compute the corresponding units dispensed during the quarter for each cell by multiplying the 5 imputed prescriptions by the unit size embedded in the NDC code. Then we multiply by the per-unit gross reimbursement price to obtain total reimbursements. We also corrected data irregularities in $609$ out of the $32{,}683$ state $\times$ NDC $\times$ quarter cells that contribute to our analysis. In these cases, the SDUD units reimbursed were positive but the corresponding Medicaid dollar amount was missing, zero, or implausible relative to per-unit price of the drug. Appendix C describes the three-tier procedure we used to resolve these cases.

\paragraph{MDRP Rebates.} The SDUD reimbursement data reflect gross payments to pharmacies at the point of dispensing. They do not include information about the quarterly rebates that manufacturers pay the states through the Medicaid Drug Rebate Program (MDRP). Under the MDRP, the per-unit rebate for \emph{brand-name drugs} depends on three manufacturer-reported prices---current quarter average manufacturer price (AMP), current quarter best price (BP), and launch quarter AMP---and equals the sum of a \emph{basic rebate} and an \emph{inflationary rebate}:
%
\begin{align*}
R_{dt}^{b} &= \underbrace{\max\bc{\bp{0.231 \times AMP_{dt}},\bp{AMP_{dt} - BP_{dt}}}}_{\text{basic rebate}} + \underbrace{\max\bc{0,\; AMP_{dt} - AMP_{dl} \times \frac{CPI_t}{CPI_l}}}_{\text{inflationary rebate}}
\end{align*}
%
The total rebate for a given NDC in a state-quarter is the per-unit rebate multiplied by units dispensed; net Medicaid spending is gross reimbursement minus the total rebate.\footnote{For generic drugs, the basic rebate is 13\% of AMP, with no best-price comparison: $R_{dt}^{g} = 0.13 \times AMP_{dt} + \max\bc{0, AMP_{dt} - AMP_{dl} \times \frac{CPI_t}{CPI_{l}}}$. Before 2024-Q1, total rebates were capped at 100 percent of AMP; the American Rescue Plan removed this cap.}  

Although the rebate formula is public, its key inputs (AMP, BP, and launch AMP) are confidential. We use the National Average Drug Acquisition Cost (NADAC) to proxy for the current and launch-quarter AMP, and the Big 4 price from the VA Federal Supply Schedule as proxy for the manufacturer's best price. Substituting the proxies and the CPI-U into the rebate formula gives an estimate of the per-unit rebate for each NDC code in each quarter; we multiply by units dispensed to obtain the total rebate for each NDC code in each state and quarter.\footnote{Appendix B describes our method for computing the rebates in more detail.} The NADAC rate is based on invoice costs from retail pharmacies. These costs include wholesale markups, so NADAC should probably be viewed as an upper bound on AMP. In contrast, the Big 4 price reflects deep federal discounts that are likely below the manufacturer's true best price. If these relationships hold, our method likely overstates rebate somewhat and yields a lower bound on net spending.

\paragraph{Compounded GLP-1 Spending.} We supplement the SDUD with debit and credit card transaction data from the Consumer Edge Brand Tracker, accessed via Dewey Data \citep{consumeredge2025brandtracker}. Consumer Edge aggregates credit and debit card transactions from a panel of U.S.\ cardholders, reporting spending and transaction counts for over 12,000 brands by state and day. We searched for brands in four categories relevant to the GLP-1 market. The first group, which we label \emph{online GLP-1 compounding platforms}, consists of six companies that have become major distributors of compounded GLP-1 medications: Hims \& Hers, Roman, Lemonaid Health, Mochi Health, Henry Meds, and Medvi. These platforms expanded rapidly after the FDA placed semaglutide on its drug shortage list beginning in 2022, which permitted compounding pharmacies to legally sell on-patent drugs.\footnote{The FDA resolved the tirzepatide shortage in October 2024 and the semaglutide shortage in February 2025, after which compounding of these products was required to wind down.} The second group---\emph{online GLP-1 prescribing platforms}---consists of nine telehealth companies that primarily facilitate prescriptions for FDA-approved GLP-1 medications: Push Health, Sequence, Form Health, Hone Health, PlushCare, Sesame Health, LifeMD, Nurx, and Peter MD. We combine the compounding and prescribing platforms into a single \emph{online GLP-1} spending measure for our event study analysis.\footnote{Noom, WeightWatchers, Found, and Calibrate also have GLP-1 prescribing arms but are classified as weight-loss platforms because their core business model centers on diet and weight management subscriptions.} These companies also sell other health and wellness products; our spending measure captures all card transactions with each merchant, not just GLP-1 purchases. The third group---\emph{weight-loss platforms}---consists of spending at eight non-GLP-1 subscription-based diet and weight management companies, including Noom, WeightWatchers, Nutrisystem, Medifast, Jenny Craig, Found, Calibrate, and Rugiet Health. The fourth is the manufacturer direct-to-consumer channel, Lilly Direct.\footnote{We also searched for additional firms from the telehealth and compounding space: Empower, Hallandale, Strive, Olympia, BelMar, Ro, Sesame Care, ShedRx, and Ivim Health. But these brands were not represented in the Consumer Edge data.} We aggregate daily brand-level transactions to the state-quarter level and scale spending per 100 Medicaid enrollee-months for comparability with the SDUD outcomes. Importantly, the Consumer Edge data clearly does not capture all private out-of-pocket spending on compounded or non-compounded GLP-1-related products. Thus, the level outcomes are probably less interesting than the trends and patterns across states.
